\section{Conclusion}

This paper asked a deliberately narrow question: under matched budget, can lightweight evidence guide Triton schedule selection better than default settings or equally budgeted naive baselines? The answer is yes, but only in a constrained and carefully evidenced sense.

The study's strongest contribution is not a large final speedup. It is a clearer account of what kinds of evidence remain useful once the budget is fixed and the workloads are realistic. Compile-adjacent signals prune better than they rank. Aligned workloads can make a selector look stronger than it is. Profiling is mixed rather than uniformly beneficial, especially once LayerNorm is written as a regime-split secondary result. Transfer failures matter because they justify retiring unstable knob families instead of stretching the narrative around them.

Within that disciplined framing, the final non-\texttt{split\_k} mainline selector delivers a small but credible representative-GEMM improvement and recovers the strongest mainline result in the project. The paper's final claim is therefore bounded on purpose: Triton tuning is not solved here, but matched-budget, bottleneck-aware tuning does yield interpretable and practically useful guidance when the evidence hierarchy and workload program are designed carefully.
